\chapter{Negative-space audit: забытые bundle-классы и непроверенный full-field carrier action}

Предыдущий corpus audit искал сохранившиеся положительные механизмы. Здесь
ставится другой вопрос: какие отрицательные вердикты были доказаны только в
узком классе и какие варианты, прямо упомянутые в ранних файлах, не получили
самостоятельного gate.

\section{Область действия прежних no-go результатов}

Сводная проверка показывает четыре важных ограничения:
\begin{enumerate}
 \item product-density no-go относится к прямым metric products;
 \item \(S^4\)-carrier ordering получен для положительного scalar heat
 state, а не для полного quantum field content;
 \item parity-even spectral-action no-go для CP не исключает новый chiral
 measure, но такой measure отсутствует в текущем parent action;
 \item конкретный \(Q_{\rm cycle}\)-vertex провален, однако общий
 single-vertex no-go в корпусе прямо помечен как недоказанный.
\end{enumerate}
Следовательно, фраза ``все варианты закрыты'' была бы переобобщением.

\section{Забытое нетривиальное circle bundle}

Файл
\path{RPFT-main/ai-promts/First-principles-00-variational.md} требует
сравнить два класса, поскольку
\begin{equation}
 H^2(\mathbb{RP}^3;\mathbb Z)\cong\mathbb Z_2.
\end{equation}
Тривиальный класс есть \(\mathbb{RP}^3\times S^1\). Homogeneous модель
нетривиального класса можно записать как
\begin{equation}
 P=\frac{S^3\times S^1}
 {(x,\theta)\sim(-x,\theta+\pi)}
 \simeq U(2).
 \label{eq:nontrivial-circle-bundle-u2}
\end{equation}
Это уже не независимый quotient базы и круга: deck transformation действует
диагонально.

Для scalar product eigenmodes
\begin{equation}
 Y_{\ell}(x)e^{im\theta}
\end{equation}
инвариантность даёт exact selection rule
\begin{equation}
 \boxed{\ell+m\equiv0\pmod2}.
 \label{eq:bundle-parity-momentum-lock}
\end{equation}
В trivial product остаются \(\ell\) even и все \(m\). Таким образом,
нетривиальное bundle действительно реализует забытый механизм:
projective parity и KK momentum больше нельзя назначать независимо.

При \(r_3=R_3/\sigma\), \(r_1=R_1/\sigma\) обе ветви имеют volume
\begin{equation}
 \frac{V}{\sigma^4}=2\pi^3r_3^3r_1.
\end{equation}
Их scalar correlation-cell densities равны
\begin{equation}
 \mathfrak f=-\frac{\log Z}{2\pi^3r_3^3r_1}.
\end{equation}

Искусственное equal-scale сечение имеет minima
\begin{align}
 r_*^{\rm product}&=0.674445979995\ldots,
 &\mathfrak f_*^{\rm product}&=-0.015642033532\ldots,\\
 r_*^{\rm bundle}&=1.274750182701\ldots,
 &\mathfrak f_*^{\rm bundle}&=-0.004035864428\ldots.
\end{align}
Но освобождение shape modulus закрывает обе ветви. При
\(r_3\to0\) и фиксированном \(r_1\)
\begin{align}
 \log Z_{\rm product}&\longrightarrow
 \log\sum_{m\in\mathbb Z}e^{-m^2/r_1^2}>0,\\
 \log Z_{\rm bundle}&\longrightarrow
 \log\sum_{n\in\mathbb Z}e^{-(2n)^2/r_1^2}>0.
\end{align}
Поэтому
\begin{equation}
 \mathfrak f_{\rm product},\mathfrak f_{\rm bundle}
 \longrightarrow-\infty.
\end{equation}

\begin{theorem}[Forgotten-bundle verdict]
Нетривиальное homogeneous \(S^1\)-расслоение является реальным пропущенным
carrier и выводит parity--momentum locking. Однако scalar correlation
density не имеет finite joint minimum и коллапсирует при сжатии базы.
Следовательно, эта ветвь закрыта как самостоятельный scalar carrier, но её
selection rule сохраняется как возможный ingredient будущего full-field
action.
\end{theorem}

\section{Что прежнее закрытие bundle не проверяло}

В S2T--II нетривиальное расслоение ранее закрывалось как rescue для C6 и
discrete selector: total space имеет continuous Wilson modulus и не даёт
нужного isolated coefficient. Это не было вычислением его normalized
carrier free energy. Настоящий gate заполняет именно этот пробел и получает
независимый отрицательный результат.

Allowed warped metrics остаются формально непросмотренными, но это пока не
конкретная теория. Произвольная warp function создаёт бесконечномерное
пространство коэффициентов. Без parent-derived cross term её оптимизация
эквивалентна введению нового функционального входа.

\begin{nogocriterion}[Запрет спасения словом ``warped'']
Warped или non-product geometry считается новым кандидатом только после
объявления action, symmetry, boundary conditions и finite parameter class.
Само снятие product ansatz не является механизмом стабилизации.
\end{nogocriterion}

\section{Главный действительно пропущенный gate}

Все последние \(S^4\) comparisons использовали scalar Gaussian operator.
Но минимальный parent ledger Тома IV содержит:
\begin{enumerate}
 \item три физических scalar operators;
 \item massive vector и полный gauge--ghost complex;
 \item две Dirac pairs;
 \item curvature endomorphisms, zero modes и determinant phases.
\end{enumerate}
Flat-limit supertrace coefficient \(67\) проверяет multiplicities, но не
определяет finite curved-space ordering.

Полный объект имеет schematically вид
\begin{equation}
 \Gamma_{\rm fluc}[M]
 =\frac12\sum_b n_b\log\det\mathcal L_b[M]
 -\frac12\sum_f n_f\log\det\mathcal D_f^2[M]
 +\Gamma_{\rm ghost}[M].
 \label{eq:full-field-carrier-fluctuation}
\end{equation}
Это signed determinant ledger, а не positive density state. Поэтому нельзя
просто заменить scalar multiplicity на supertrace внутри прежней entropy
формулы: отрицательные fermionic weights разрушили бы positivity \(\rho\).

Кроме того, \(S^4\) и \(S^2\times S^2\) имеют разные local curvature
invariants и topology. Их divergent/local determinant pieces не обязаны
сокращаться в разности. До carrier comparison необходимо одной схемой
фиксировать cosmological, Einstein и quadratic-curvature counterterms.

\begin{proposition}[Настоящий незакрытый вариант]
Полный renormalized gauge--ghost--scalar--Dirac fluctuation action на
\(S^4\) и \(S^2\times S^2\) в корпусе не вычислен. Поэтому scalar
correlation ordering не является full-field vacuum theorem. Это не лазейка
в уже доказанном no-go, а отсутствующий следующий член TOE~6.5 functional.
\end{proposition}

\section{Другие найденные, но менее сильные пропуски}

Общий single-vertex no-go действительно не доказан. Для его проверки нужно
заранее ограничить finite algebra volume/winding vertices и exhaustive
readout class. Этот путь может локально восстановить charged-lepton seed,
но сам по себе не решает carrier, scale и CP gates, поэтому имеет второй
приоритет.

Старый quark/Skyrmion файл предлагает полноценную lattice minimization, но
его action coefficients и связь с общим parent functional постулированы.
Вычисление было бы тестом отдельной phenomenological soliton model, а не
пропущенным замыканием текущей теории.

Hopf-fibration, Leech-lattice, blackbody и string analogies в архиве
переиспользуют уже известные числа или требуют новую категорию полей. После
deduplication они не дают независимого механизма measure.

\section{Новая очерёдность}

Negative-space audit задаёт следующий порядок:
\begin{enumerate}
 \item вычислить \(\Gamma_{\rm fluc}^{\rm ren}\) из
 \eqref{eq:full-field-carrier-fluctuation} на fixed-volume
 \(S^4\) и \(S^2\times S^2\) в одной preregistered scheme;
 \item вставить эту величину в joint state--carrier functional TOE~6.5 и
 проверить полный Hessian;
 \item возвращаться к warped/bundle geometries только если cross term
 возникает из того же determinant, а не добавляется вручную;
 \item отдельно exhaust finite single-vertex algebra, не смешивая этот
 локальный вопрос с carrier selection.
\end{enumerate}

\begin{theorem}[Negative-space final verdict]
Во всём массиве найдены не только повторения. Один забытый вариант ---
нетривиальное \(S^1\)-расслоение --- теперь проверен и закрыт как scalar
vacuum, сохранив exact parity--momentum rule. Главный непроверенный вариант
--- полный renormalized field fluctuation determinant на competing
carriers. Именно он, а не очередное численное совпадение, является следующим
содержательным расчётом программы.
\end{theorem}

Машинный аудит сохранён в
\path{s2t_v4_negative_space_bundle_fluctuation_gate.py} и
\path{s2t_v4_negative_space_bundle_fluctuation_gate_results.json}.